% !TeX root = ../tesis.tex

\section{Eritrocito sano}
\label{section:Res_Eritrocito_Sano}

Para identificar las propiedades ópticas de los eritrocitos en las bandas de Soret, $\beta$ y $\alpha$, se consideró el rango espectral de 400  a 600 nm.  A lo largo de esta sección, por comodidad se define a una \textit{dirección} como un ángulo de esparcimiento específico ($\theta=0^\circ$ para el esparcimiento frontal, $\theta=90^\circ$ para el lateral y $\theta=180^\circ$ para el retroesparcimiento), mientras que una \textit{región} corresponde a un intervalo angular ($0^\circ\leq\theta<60^\circ$ para el esparcimiento frontal, $60^\circ\leq\theta<120^\circ$ para el lateral y $120^\circ\leq\theta\leq180^\circ$ para el retroesparcimiento).%
% --- JAUA ---
% te propongo mejor ennumerar los nombres y entrea paréntesis los valores. Por ejemplo:
% se definen la dirección frontal como ek pangulo de esparcimeinto de \theta =0, la lateral de 90° y la dirección de tresotespacimieto para 180°. Por otro lado, se defienen tambiénn las regiones de esparcimiento frontal (pones los valores), el lateral (cvalores) y.... 
% Creo que eso lo haría más claro de leer y ponemos peso en todos los nombres que utilizas.
%
  Para modelar un eritrocito sano, se empleó un óvalo de Cassini %
  % --- JAUA ---
  % con los parámetreos mostrados en la Fig XX con....
  %
   como el de la Fig. \ref{subfig:Shape_health} con parámetros de $a=2.6\,\mu\text{m}$, $b=2.7\,\mu\text{m}$ y $c=0.812\,\mu\text{m}$, los cuales resproducen el radio equivalente de 3.75 $\mu$m y el volumen de 80.63~$\mu\text{m}^3$ reportados por Ergül et al. en la Ref.~\cite{ergulComputationalStudyScattering2010}. De acuerdo con los resultados obtenidos en la Sección~\ref{section:optical_properties}, para el plasma se consideró una función dieléctrica constante de $\varepsilon_{\text{Plasma}}=1.839$, %
   % --- JAUA ---
   % en secciones anteriores habiamos modificado a 1.84; recomiendo dejarlo de esa forma y omitir que cambia para cuatro ciffras decimales.
   %
    ya que su parte imaginaria resulta despreciable y la parte real presenta variaciones únicamente en la cuarta cifra decimal dentro del rango espectral de interés. Por su parte, para el eritrocito se empleó la función dieléctrica de una CHCM de 30.6~g/dL %
	% --- JAUA ---
	% Creo que hay que enfatizar que un eritrosito sano corresponde a esa concentración y que no es una elección arbitraria; te recomiendo hacer mención de las figuras correspodientes en secciones anterior3es, o reecoradar que esto se discutió ya antes.
	%
	. El eritrocito se orientó con su eje de simetría a lo largo de la dirección $z$ y se iluminó con una onda plana que se propaga en la misma dirección, polarizada en $x$.

En la Fig. \ref{subfig:Qext_health} se muestra la eficiencia de extinción $Q_{\text{ext}}$ en función de la longitud de onda $\lambda$ de la onda incidente; se resaltan en gris las bandas que rodean los mínimos asociados a la banda de Soret y a las bandas $\beta$ y $\alpha$ localizados en 431 nm, 541 nm y 573 nm, respectivamente, los cuales se consideran longitudes de onda representativas para el análisis de los patrones de esparcimiento. Se observa que la respuesta espectral del eritrocito sigue la dependencia espectral de su función dieléctrica, reflejando la influencia de las transiciones ópticas de la hemoglobina como mínimos en $Q_{\text{ext}}$.
%
\begin{figure}[h!]
	\centering
	\sidesubfloat[First image]{\hspace{-0.3cm}{
			\includegraphics[width=0.45\textwidth]{4-Resultados/Eritrocito_sano_corte.pdf}\label{subfig:Shape_health}}}\hspace{0.1cm}
	\sidesubfloat[Second image]{\hspace{-0.5cm}{	\includegraphics[width=0.45\textwidth]{4-Resultados/Qext_EriSano.pdf}\label{subfig:Qext_health}}}
	\vspace{-0.1cm}
	\caption{\textbf{a)} Modelo de eritrocito sano basado en óvalos de Cassini con parámetros de $a=2.6\,\mu\text{m}$, $b=2.7\,\mu\text{m}$ y 
		$c=0.812\,\mu\text{m}$, que corresponden a un radio de $3.75\,\mu$m. \textbf{b)} Eficiencia de extinción $Q_{\text{ext}}$ en función de la longitud de onda $\lambda$ de la onda incidente para un eritrocito orientado con su eje de simetría en el eje $z$ iluminado con una onda en dirección del eje $z$ y polarizada en $x$. La línea continua que une los puntos es una guía visual. Se resaltan en gris las bandas que 
		rodean los mínimos asociados a la banda de Soret y a las bandas Q.}
	\label{fig:Shape_Qext_health}
\end{figure}
%

En esta sección se discute únicamente el patrón de esparcimiento de $\lambda=431$~nm asociado a la banda de Soret, mientras que los de $\lambda=541$ y $573$ nm se incluyen en el Apéndice \ref{section:apendix3}. En la Fig.~\ref{subfig:FF_431_health} se muestra la sección transversal diferencial de esparcimiento $\dd{C_{\text{sca}}}/\dd{\Omega}$ para $\lambda=431$~nm. La magnitud de $\dd{C_{\text{sca}}}/\dd{\Omega}$ se encuentra normalizada respecto a su valor máximo, que corresponde al valor más saturado. Se observa que $\dd{C_{\text{sca}}}/\dd{\Omega}$ es mayor en la dirección frontal de propagación, lo cual se refleja en un factor de anisotropía\footnote{%
	El factor de anisotropía $g$ describe la direccionalidad de la luz esparcida por una partícula y corresponde al promedio ponderado de la intensidad de esparcimiento normalizada con el coseno del ángulo de esparcimiento \cite{videenAsymmetryParameterAggregate1998}. Se define como \cite{bohrenAbsorptionScatteringLight2008}
	$$g=\dfrac{1}{C_{\text{sca}}}\int_{4 \pi}\dfrac{\dd{C_{\text{sca}}}}{\dd{\Omega}}\cos\theta\dd{\Omega},$$
	por lo que $g$ varía entre $-1$ y 1. En particular, para una partícula que esparce la luz de forma isótropa o si su esparcimiento es simétrico respecto a un ángulo de esparcimiento $\theta=90^{\circ}$, $g$ se anula. Si la partícula esparce más luz hacia la dirección frontal ($\theta=0^{\circ}$), $g>0$; si el esparcimiento se dirige más hacia la dirección posterior ($\theta=180^{\circ}$), $g<0$ \cite{bohrenAbsorptionScatteringLight2008}.} %
$g$ igual a 0.97. Se resaltan además, dos planos de observación: $\varphi=0^{\circ}$ (azul claro), que es paralelo a la polarización de la onda indicente y $\varphi=90^{\circ}$ (azul oscuro), perpendicular a la polarización de la onda incidente. Con el fin de estudiar el efecto de la polarización de la onda incidente en el esparcimiento, en la Fig. \ref{subfig:FF_polar_431_health} se muestra $\dd{C_{\text{sca}}}/\dd{\Omega}$ en decibeles en función del ángulo de esparcimiento $\theta$ para dos cortes correspondientes a los planos $\varphi=0^{\circ}$ (línea continua azul claro) y $\varphi=90^{\circ}$ (línea continua azul oscuro). En la Fig. \ref{subfig:FF_cartesian_431_health} se muestra $\dd{C_{\text{sca}}}/\dd{\Omega}$ en función del ángulo de esparcimiento para ambos cortes en esala logarítimica. Debido a que las dimensiones del eritrocito son mayores que la longitud de onda incidente, su respuesta se encuentra fuera del régimen dipolar. En consecuencia, el patrón de esparcimiento presenta una estructura angular compleja, asociada con los retardos de fase y la interferencia entre las contribuciones al campo esparcido provenientes de distintas regiones de la célula.
	%
	\begin{figure}[t]
		\centering
		\sidesubfloat[First image]{\hspace{-0.5cm}{	\includegraphics[width=7.3cm]{4-Resultados/FFdistribution431_EriSano.pdf}\label{subfig:FF_431_health}}}\hspace{0.1cm}
		\sidesubfloat[Second image]{\hspace{-0.5cm}{
				\includegraphics[width=8cm]{4-Resultados/polarplot431FF_EriSano.pdf}\label{subfig:FF_polar_431_health}}}\vspace{-0.1cm}
		\
		\vspace{0.1cm}
		
		\sidesubfloat[Third image]{\hspace{0.3cm}{
				\includegraphics[width=14cm]{4-Resultados/linearplot431FF_EriSano.pdf}\label{subfig:FF_cartesian_431_health}}}\vspace{-0.1cm}
		\caption{\textbf{a)} Sección transversal diferencial de esparcimiento $\dd{C_{\text{sca}}}/\dd{\Omega}$ normalizada respecto a su valor máximo para una longitud de onda $\lambda=431$ nm en función del ángulo de esparcimiento $\theta$ y del ángulo $\varphi$ para un eritrocito sano. Se identifican dos planos $\varphi=0^{\circ}$ en azul claro y $\varphi=90^{\circ}$ en azul oscuro. \textbf{b)} Sección transversal diferencial de esparcimiento $\dd{C_{\text{sca}}}/\dd{\Omega}$ en decibeles en función de $\theta$ para dos cortes: $\varphi=0^{\circ}$ (puntos azules claro) y $\varphi=90^{\circ}$ (puntos azules oscuro). \textbf{c)} Sección transversal diferencial de esparcimiento $\dd{C_{\text{sca}}}/\dd{\Omega}$ en función de $\theta$ para dos cortes: $\varphi=0^{\circ}$ (puntos azules claro) y $\varphi=90^{\circ}$ (puntos azules oscuro). La línea continua que une los puntos es una guía visual.}
		\label{fig:FF_431_health}
	\end{figure}
	%
	
Para cuantificar tanto la distribución angular como la magnitud de esparcimiento en cada región, se calculó el coeficiente de correlación de Pearson\footnote{
	El coeficiente de correlación de Pearson entre dos conjuntos de datos $x$ y $y$ se define como \cite{rossIntroductionProbabilityStatistics2014}
	$$r=\dfrac{\sum_{i}{(x_i-\overline{x})(y_i-\overline{y})}}{\sqrt{\sum_{i}{(x_i-\overline{x})^2} \sum_{i}{(y_i-\overline{y})^2}}},$$
	donde $\overline{x}$ indica el promedio del conjunto de datos. Por definición, $r$ es un valor entre $-1$ y 1, donde su valor absoluto es una medida de la fuerza de la relación lineal entre los dos pares de datos, mientras que su signo indica la dirección de la relación lineal \cite{rossIntroductionProbabilityStatistics2014}.} % 
	%
	% --- JAUA ---
	% Te adelanto que alo mejor vas a tener que explicar que no significa que sean iguales, pero que siguen mismas tendencias (la ecuación te puede aydar en donde además puedes aprovechar la dependencia en x. Creo que en general pensamos en mínimos cuaddrrados donde eso no ocurre.)
	%
entre ambos planos de observación y se integró $\dd{C_{\text{sca}}}/\dd{\Omega}$ en los intervalos correspondientes al esparcimiento frontal, lateral y de retroesparcimiento. Ambas cantidades proporcionan información complementaria: mientras que el coeficiente de Pearson permite comparar la estructura angular de los patrones, independientemente de diferencias globales de amplitud, la integral de la sección transversal diferencial permite cuantificar el campo esparcido. En la región frontal se obtiene $r=1$, lo que indica que las posiciones y amplitudes relativas de los máximos y mínimos producidos por interferencia presentan la misma dependencia angular en ambos planos de observación. Asimismo, las contribuciones integradas normalizadas por el área $S=4.41\times 10^{7}\,\text{nm}^2$ son $1,817.17\times 10^{-4}$ y $1,788.52\times 10^{-4}$ para $\varphi=0^{\circ}$ y $\varphi=90^{\circ}$, respectivamente, por lo que no solo la estructura angular, sino también la magnitud del esparcimiento frontal, presenta una baja dependencia con el plano de observación.  Este mismo caso ocurre en la región de retroesparcimiento, donde $r=0.99$ y las contribuciones integradas son iguales a $2.01\times 10^{-4}$ y $2.97\times 10^{-4}$ para $\varphi=0^{\circ}$ y $\varphi=90^{\circ}$, respectivamente. En contraste, en la región lateral $r= 0.67$ y las contribuciones integradas son de $11.71\times 10^{-4}$ y $21.70\times 10^{-4}$ para $\varphi=0^{\circ}$ y $\varphi=90^{\circ}$, respectivamente. En esta región, las diferencias tanto en la estructura de máximos y mínimos como en la magnitud integrada indican una mayor dependencia del patrón de esparcimiento con el plano de observación asociado a la polarización lineal. A partir de estos resultados, se observa que la contribución frontal supera hasta dos órdenes de magnitud a las contribuciones lateral y de retroesparcimiento, lo que muestra la direccionalidad del esparcimiento del eritrocito. Este comportamiento puede asociarse con que el campo esparcido por el eritrocito resulta de la superposición de las contribuciones generadas por la distribución de la polarización inducida en toda la célula y, debido a que sus dimensiones son mayores que la longitud de onda incidente, existen retardos de fase apreciables entre las contribuciones provenientes de distintas regiones del eritrocito. En la dirección frontal, las diferencias de fase son menores, favoreciendo una superposición predominantemente constructiva, mientras que al aumentar el ángulo de esparcimiento, las diferencias de fase entre las distintas contribuciones se incrementan, produciendo mayor interfencia destructiva y en consecuencia, la reducción del esparcimiento lateral y posterior.
	
	
	

En la Tabla~\ref{table:val_eri} se resumen, para las tres longitudes de onda representativas, los valores del factor de anisotropía $g$, la correlación en la región de esparcimiento lateral entre los cortes $\varphi=0^\circ$ y $\varphi=90^\circ$, así como las contribuciones de la sección transversal de esparcimiento integradas en las regiones frontal, lateral y de retroesparcimiento para ambos cortes. La correlación se presenta únicamente para la región lateral debido a que, en las regiones frontal y de retroesparcimiento, los patrones permanecen invariantes entre ambos cortes, con coeficientes de correlación de 1 y 0.99, respectivamente. Se observa que el factor de anisotropía aumenta de 0.97 a 0.99 al pasar de la banda de Soret a la banda $\alpha$, lo que es consistente con el predominio del esparcimiento frontal en todo el intervalo espectral analizado. Asimismo, la contribución integrada en la región frontal aumenta con la longitud de onda, mientras que las regiones lateral y de retroesparcimiento continúan aportando una fracción menor al esparcimiento total. Por otra parte, la correlación en la región lateral no presenta una tendencia creciente con la longitud de onda, alcanzando su valor mínimo en la banda $\beta$ ($r=0.49$), lo que indica que en esta región espectral la distribución angular del esparcimiento es la más sensible al plano de observación. Sin embargo, la diferencia entre las contribuciones integradas entre los dos planos es inferior (38.93 \%) para esta misma longitud de onda que para $\lambda=431$~nm (85.31 \%) y $\lambda=573$~nm (98.35 \%). Esto evidencia que una menor correlación entre dos planos de observación no implica necesariamente una mayor diferencia en su contribución en $C_{\text{sca}}$.

%
\begin{table}[h!]
	\begin{tabular}{ccccccccc}
		\hline
		\multirow{3}{*}{$\lambda$ [nm]} & \multirow{3}{*}{$g$} & \multirow{3}{5em}{Correlación lateral} & \multicolumn{6}{c}{$C_{\text{sca}}/S\times 10^{-4}$ } \\
		\cline{4-9}
		& &  &   \multicolumn{2}{c}{Frontal} & \multicolumn{2}{c}{Lateral} & \multicolumn{2}{c}{Posterior}\\
		\cline{4-9}
		&  &  &  $\varphi=0^{\circ}$ & $\varphi=90^{\circ}$ &$\varphi=0^{\circ}$ & $\varphi=90^{\circ}$ &$\varphi=0^{\circ}$ & $\varphi=90^{\circ}$ \\ \hline
		431 (Soret) & 0.97  & 0.67 & 1,817.17 & 1,788.52 & 11.71 & 21.70 & 2.01 & 2.97  \\
		541 ($\beta$) & 0.98 & 0.49 & 4,324.20 & 4,296.07 & 15.08 & 20.95 & 10.78 & 13.77\\
		573 ($\alpha$) & 0.99 & 0.71 & 4,650.05 & 4,626.42 & 12.75 & 25.29 & 2.46 & 3.41\\
		\hline
	\end{tabular}
	\caption{Valores del factor de anisotropía $g$, la correlación en la región de esparcimiento lateral entre los cortes $\varphi=0^\circ$ y $\varphi=90^\circ$ y las contribuciones de la sección transversal de esparcimiento integradas en las regiones frontal, lateral y de retroesparcimiento para $\varphi=0^{\circ}$ y $\varphi=90^{\circ}$ para $\lambda=431,541$ y 573 nm. }
	\label{table:val_eri}
	\vspace{0.2cm} 
\end{table}
%




